\chapter{Objetivos concretos y metodología de trabajo}
\section{Objetivo general}
El objetivo general de este trabajo es desarrollar y validar un sistema de predicción de resultados de partidos de fútbol basado en técnicas de \textit{machine learning}, capaz de estimar la probabilidad de victoria local, empate y victoria visitante a partir de datos de eventos transformados en métricas segmentadas temporalmente. El proyecto incluye la adquisición y preparación de los datos, el entrenamiento del modelo, la comparación con distintos algoritmos y la validación de su rendimiento, estableciendo como criterio de éxito alcanzar o aproximarse a una precisión del 70\,\% sobre un conjunto de test, complementada con métricas adicionales que permitan valorar el rendimiento del sistema.

\section{Objetivos específicos}
Con el fin de que el objetivo general resulte específico, medible, alcanzable, relevante y acotado al marco temporal del proyecto, se plantean los siguientes objetivos específicos:

\begin{itemize}
    \item Realizar, durante la fase inicial del proyecto, una revisión bibliográfica sobre estudios existentes de predicción de resultados deportivos. El propósito de este objetivo es extraer criterios claros sobre tipos de datos, variables, algoritmos y métricas de evaluación que sirvan de base para el diseño experimental posterior.
    \item Recopilar y tratar los datos procedentes de \textit{StatsBomb Open Data}, llegando a obtener un repositorio de datos consistente y analizando sus limitaciones y ventajas para la predicción del resultado final.
    \item Construir un conjunto de variables explicativas agrupadas por segmentos temporales que describan el estado y la evolución del partido, incluyendo métricas acumuladas de rendimiento de ambos equipos. Este objetivo será satisfactorio si las variables generadas permiten representar de forma suficientemente informativa la evolución del encuentro.
    \item Entrenar y comparar varios modelos de clasificación multiclase adecuados para el problema, incluyendo al menos modelos \textit{ensemble} y, si la disponibilidad de datos y tiempo lo permite, algún enfoque basado en aprendizaje profundo. Posteriormente, se seleccionarán los modelos con mejores rendimientos predictivos y se analizarán sus resultados en un conjunto de datos de test.
    \item Evaluar el sistema obtenido mediante métricas cuantitativas apropiadas para clasificación multiclase, incluyendo como mínimo \textit{accuracy} y medidas complementarias como \textit{F1-score}, y analizar su rendimiento en función del segmento temporal de juego. De forma específica, se establecerá como referencia de éxito alcanzar o aproximarse de manera justificada al 70\,\% de precisión en test, interpretando además en qué contextos temporales la predicción resulta más fiable.
    \item Integrar el modelo final en una solución que permita introducir o simular la evolución de un partido y visualizar la probabilidad estimada de cada resultado a lo largo del tiempo. Este objetivo se dará por cumplido cuando el sistema final pueda utilizarse como prueba de concepto funcional del enfoque desarrollado en el proyecto.
\end{itemize}

\section{Metodología del trabajo}
Para alcanzar los objetivos planteados en este proyecto, se adopta una metodología ágil que combina dos marcos de trabajo complementarios: \textbf{Scrum} \citep{schwaber2020scrum}, que se encarga de la gestión del proyecto, y \textbf{CRISP-DM} (\textit{Cross Industry Standard Process for Data Mining}) \citep{chapman2000crispdm}, que se enfoca en el ciclo de vida técnico de un proyecto basado en datos. El uso conjunto de ambos marcos permite adaptar la planificación del proyecto a los resultados que se van obteniendo en cada fase.

\textbf{Scrum} se emplea para estructurar el desarrollo del proyecto en iteraciones cortas, de modo que el alcance y las prioridades puedan ajustarse en cada ciclo a partir de los resultados obtenidos. Este enfoque es especialmente adecuado en contextos donde la validez de las decisiones, como la selección de variables o la elección del modelo, solo puede comprobarse mediante la experimentación. Además, este marco de trabajo incorpora eventos específicos orientados a la planificación, la evaluación del progreso y la mejora continua del proceso.

Por otro lado, \textbf{CRISP-DM} es un marco ampliamente utilizado para proyectos de minería y análisis de datos y, por su estructura y metodología, puede ser adaptado también para proyectos de inteligencia artificial como el planteado en este trabajo. CRISP-DM aborda el desarrollo en seis fases que se revisan y ajustan a lo largo del proyecto: comprensión del problema, comprensión de los datos, preparación de los datos, modelado, evaluación y despliegue. La naturaleza cíclica de este marco es especialmente relevante en el contexto del proyecto, ya que los resultados obtenidos en fases avanzadas pueden llevar a revisar decisiones tomadas anteriormente, como la selección de variables o el preprocesamiento de los datos.

\subsection{Fases del proyecto}
Las seis fases de CRISP-DM se particularizan para el presente trabajo de la siguiente manera:

\subsubsection{Fase 1. Comprensión del problema (\textit{Business Understanding})}
En esta primera fase se define el alcance del proyecto y se establecen las condiciones que determinan su éxito. El objetivo es desarrollar un sistema que pueda predecir la probabilidad de cada resultado posible en un partido de fútbol, basándose en métricas acumuladas durante el partido. Se identifican los principales ámbitos de aplicación del sistema, como el análisis táctico por parte del cuerpo técnico, la generación de contenidos para \textit{broadcasting} y los sistemas de apuestas deportivas \citep{goes2021bigdata, rein2016bigdata}. A partir de este análisis, el objetivo se formaliza como una tarea de clasificación multi-clase, con un criterio de éxito de alcanzar una precisión superior al 70\,\% sobre un conjunto de test.

\subsubsection{Fase 2. Comprensión de los datos (\textit{Data Understanding})}
La fuente de datos principal del proyecto es el repositorio \textit{StatsBomb Open Data} \citep{statsbomb2024opendata}, que proporciona registros detallados de eventos de partidos de diversas competiciones y temporadas. En esta fase se analizan los datos disponibles para caracterizar su estructura, evaluar su calidad y detectar posibles limitaciones. Se lleva a cabo un análisis exploratorio de datos que orienta las decisiones de las fases posteriores; su desarrollo completo se recoge en el capítulo de desarrollo.

\subsubsection{Fase 3. Preparación de los datos (\textit{Data Preparation})}
Esta fase tiene como objetivo transformar los datos brutos del repositorio \textit{StatsBomb} en un dataset estructurado y listo para el modelado. A partir de los eventos individuales se construye un dataset tabular con métricas acumuladas por segmentos temporales, al que se aplican técnicas de limpieza e ingeniería de características \citep{kuhn2019feature}. El detalle de las variables generadas y las transformaciones aplicadas se describe en el capítulo de desarrollo.

\subsubsection{Fase 4. Modelado (\textit{Modeling})}
En esta fase se entrenan y comparan distintos algoritmos de clasificación multi-clase con el objetivo de evaluar diferentes estrategias de aprendizaje sobre el mismo conjunto de datos y bajo un procedimiento experimental idéntico. Para establecer un punto de comparación, se toma la predicción sistemática de la clase mayoritaria (victoria local), que alcanza en torno al 47\,\% de precisión, por lo que cualquier modelo propuesto debe superarla con claridad. Se contemplan métodos \textit{ensemble} y un clasificador lineal como alternativa.

El conjunto de datos se divide de forma temporal por temporada. El 80\,\% de los partidos más antiguos se destina al entrenamiento y el 20\,\% más reciente al test. La división se realiza a nivel de partido para evitar que minutos de un mismo encuentro aparezcan a la vez en entrenamiento y en test. Dentro del conjunto de entrenamiento, cada modelo utiliza una partición interna de validación para la selección de hiperparámetros, respetando también el orden cronológico. Para controlar el sobreajuste se aplican pesos de clase balanceados y parámetros de regularización propios de cada algoritmo, y se comprueba que la divergencia entre las métricas de entrenamiento y validación se mantiene dentro de márgenes razonables.

\subsubsection{Fase 5. Evaluación (\textit{Evaluation})}
Los modelos entrenados en la fase anterior se comparan utilizando un conjunto de test formado por datos de partidos no utilizados durante el entrenamiento y un conjunto de métricas específicas \citep{sokolova2009metrics}. Además, se analiza la fiabilidad del resultado según el momento del partido en que se realiza la predicción, ya que de esta forma se puede determinar a partir de qué instante del partido las predicciones resultan más fiables.

\subsubsection{Fase 6. Despliegue (\textit{Deployment})}
El modelo con mejores métricas de predicción se integra en una plataforma interactiva que permite introducir o simular datos de un partido y visualizar la evolución de las probabilidades de cada resultado minuto a minuto.

\subsection{Planificación y cronograma}

El desarrollo del proyecto se ha estructurado en sprints de dos semanas siguiendo el marco Scrum. La tabla~\ref{tab:cronograma} recoge la planificación ejecutada, indicando para cada sprint el período, el objetivo principal y las actividades completadas.

\begin{table}[H]
    \centering
    \begin{tabularx}{\textwidth}{|c|p{3cm}|X|}
        \hline
        \textbf{Sprint} & \textbf{Período}       & \textbf{Objetivo y actividades principales}                                                                                                \\
        \hline
        1               & \textit{25/03 - 08/04} & Revisión bibliográfica y selección de fuentes de datos. Estudio del repositorio StatsBomb Open Data y definición del alcance del problema. \\
        2               & \textit{01/04 - 17/04} & Descarga y construcción del dataset de eventos. Implementación del script de descarga y verificación de integridad.                        \\
        3               & \textit{17/04 - 01/05} & Análisis exploratorio de datos (EDA). Estudio de la distribución de clases, análisis temporal, correlaciones y limitaciones del dataset.   \\
        4               & \textit{01/05 - 15/05} & Ingeniería de características. Construcción del dataset de modelado con variables acumuladas, de ventana deslizante y derivadas.           \\
        5               & \textit{15/05 - 29/05} & Modelado inicial. Entrenamiento del modelo Random Forest con búsqueda de hiperparámetros.                                                  \\
        6               & \textit{29/05 - 12/06} & Modelado comparativo. Entrenamiento de XGBoost y SVM lineal. Evaluación comparativa de los tres modelos.                                   \\
        7               & \textit{29/05 - 12/06} & Análisis de resultados e interpretabilidad. Análisis de importancia de variables y evaluación temporal por minuto.                         \\
        8               & \textit{01/04 - 24/06} & Redacción de la memoria y desarrollo del prototipo de despliegue.                                                                          \\
        \hline
    \end{tabularx}
    \caption{Cronograma ejecutado del proyecto por sprints}
    \label{tab:cronograma}
\end{table}

La naturaleza cíclica de Scrum ha resultado especialmente adecuada para este proyecto, dado que varias decisiones de modelado solo pudieron validarse mediante experimentación. Por ejemplo, la inclusión inicial de un modelo LSTM fue revisada al final del sprint de modelado inicial, cuando se constató que la complejidad adicional no estaba justificada dado el enfoque de agregación temporal adoptado. En su lugar, se incorporó un SVM lineal que permite valorar el rendimiento de una frontera de decisión simple frente a los métodos \textit{ensemble}. Esta capacidad de ajustar el alcance entre iteraciones es precisamente la ventaja que aporta Scrum en un proyecto de ciencia de datos donde los resultados experimentales condicionan las decisiones siguientes.

\subsection{Herramientas y tecnologías}
El desarrollo se realiza íntegramente en Python porque es el lenguaje de programación líder para inteligencia artificial y ciencia de datos gracias a su amplio ecosistema de librerías especializadas \citep{raschka2020python}, que cubren desde el preprocesamiento de los datos hasta el modelado y la evaluación de modelos de aprendizaje automático. Para la obtención de los datos se utiliza \texttt{statsbombpy}, que permite descargar los datos del repositorio \textit{StatsBomb} directamente desde Python. El control de versiones tanto del código como de la memoria del proyecto se gestiona con Git y GitHub porque permite mantener la trazabilidad del desarrollo a lo largo del proyecto.

\section{Grado de cumplimiento de los objetivos}

La tabla~\ref{tab:cumplimiento} recoge el grado de cumplimiento de cada objetivo específico al cierre del proyecto, indicando el estado alcanzado y las desviaciones respecto a la planificación inicial.

\begin{table}[H]
    \centering
    \begin{tabularx}{\textwidth}{|p{5cm}|p{2.8cm}|X|}
        \hline
        \textbf{Objetivo específico}                                       & \textbf{Estado}       & \textbf{Observaciones}                                                                                                                                                                                                                                                                                                                 \\
        \hline
        Revisión bibliográfica sobre predicción de resultados deportivos   & Cumplido              & Se ha ampliado con trabajos recientes sobre Transformers, redes de grafos y calibración probabilística.                                                                                                                                                                                                                                \\
        \hline
        Recopilación y tratamiento de datos de StatsBomb Open Data         & Cumplido              & Dataset de 12,2\,M de eventos descargado, verificado y almacenado en \path{datasets/statsbomb_events.csv}.                                                                                                                                                                                                                             \\
        \hline
        Construcción de variables explicativas por segmentos temporales    & Cumplido              & Dataset de modelado con 48 variables generado: acumuladas, ventana deslizante y derivadas (\path{datasets/match_minute_features.csv}).                                                                                                                                                                                                 \\
        \hline
        Entrenamiento y comparación de modelos de clasificación multiclase & Parcialmente cumplido & Se han entrenado tres modelos (Random Forest, XGBoost y SVM lineal). No se ha incluido aprendizaje profundo: el LSTM previsto fue descartado por complejidad adicional sin beneficio claro sobre el enfoque de agregación adoptado.                                                                                                    \\
        \hline
        Evaluación mediante métricas cuantitativas y análisis temporal     & Cumplido              & Los tres modelos se evalúan con \textit{accuracy}, F1-score macro y \textit{log-loss} en seis puntos temporales (min 15, 30, 45, 60, 75 y 90). El umbral del 70\,\% de precisión no se alcanza a nivel global (mejor resultado: 65,4\,\%), pero sí desde el minuto 60 en adelante, lo que es coherente con la naturaleza del problema. \\
        \hline
        Integración del modelo en una solución interactiva                 & Cumplido              & El modelo final XGBoost se ha integrado en un prototipo web desarrollado con Streamlit. La aplicación permite explorar partidos del conjunto de test, visualizar la evolución de probabilidades y simular estados de partido.                                                                                                          \\
        \hline
    \end{tabularx}
    \caption{Grado de cumplimiento de los objetivos específicos del proyecto}
    \label{tab:cumplimiento}
\end{table}

Las principales desviaciones respecto a la planificación inicial son dos. En primer lugar, la exclusión del LSTM como alternativa de modelado: aunque se planteó como posibilidad condicionada a la disponibilidad de datos y tiempo, la experimentación mostró que el enfoque de agregación temporal minuto a minuto no requiere modelar la secuencia completa de eventos para obtener resultados competitivos. En segundo lugar, el objetivo del 70\,\% de precisión global no se ha alcanzado con ninguno de los modelos; no obstante, el análisis temporal muestra que esta limitación es característica de los primeros minutos del partido, donde la información disponible es escasa, y que los modelos superan ese umbral a medida que avanza el encuentro. Ambas desviaciones están justificadas por los resultados obtenidos y se discuten en detalle en los capítulos de Desarrollo y Resultados.
